% Options for packages loaded elsewhere
\PassOptionsToPackage{unicode}{hyperref}
\PassOptionsToPackage{hyphens}{url}
%
\documentclass[
]{article}
\usepackage{lmodern}
\usepackage{amssymb,amsmath}
\usepackage{ifxetex,ifluatex}
\ifnum 0\ifxetex 1\fi\ifluatex 1\fi=0 % if pdftex
  \usepackage[T1]{fontenc}
  \usepackage[utf8]{inputenc}
  \usepackage{textcomp} % provide euro and other symbols
\else % if luatex or xetex
  \usepackage{unicode-math}
  \defaultfontfeatures{Scale=MatchLowercase}
  \defaultfontfeatures[\rmfamily]{Ligatures=TeX,Scale=1}
\fi
% Use upquote if available, for straight quotes in verbatim environments
\IfFileExists{upquote.sty}{\usepackage{upquote}}{}
\IfFileExists{microtype.sty}{% use microtype if available
  \usepackage[]{microtype}
  \UseMicrotypeSet[protrusion]{basicmath} % disable protrusion for tt fonts
}{}
\makeatletter
\@ifundefined{KOMAClassName}{% if non-KOMA class
  \IfFileExists{parskip.sty}{%
    \usepackage{parskip}
  }{% else
    \setlength{\parindent}{0pt}
    \setlength{\parskip}{6pt plus 2pt minus 1pt}}
}{% if KOMA class
  \KOMAoptions{parskip=half}}
\makeatother
\usepackage{xcolor}
\IfFileExists{xurl.sty}{\usepackage{xurl}}{} % add URL line breaks if available
\IfFileExists{bookmark.sty}{\usepackage{bookmark}}{\usepackage{hyperref}}
\hypersetup{
  pdftitle={Network Analysis on Rumination},
  hidelinks,
  pdfcreator={LaTeX via pandoc}}
\urlstyle{same} % disable monospaced font for URLs
\usepackage[margin=1in]{geometry}
\usepackage{graphicx,grffile}
\makeatletter
\def\maxwidth{\ifdim\Gin@nat@width>\linewidth\linewidth\else\Gin@nat@width\fi}
\def\maxheight{\ifdim\Gin@nat@height>\textheight\textheight\else\Gin@nat@height\fi}
\makeatother
% Scale images if necessary, so that they will not overflow the page
% margins by default, and it is still possible to overwrite the defaults
% using explicit options in \includegraphics[width, height, ...]{}
\setkeys{Gin}{width=\maxwidth,height=\maxheight,keepaspectratio}
% Set default figure placement to htbp
\makeatletter
\def\fps@figure{htbp}
\makeatother
\setlength{\emergencystretch}{3em} % prevent overfull lines
\providecommand{\tightlist}{%
  \setlength{\itemsep}{0pt}\setlength{\parskip}{0pt}}
\setcounter{secnumdepth}{-\maxdimen} % remove section numbering

\title{Network Analysis on Rumination}
\author{}
\date{\vspace{-2.5em}}

\begin{document}
\maketitle

{
\setcounter{tocdepth}{2}
\tableofcontents
}
\hypertarget{methods}{%
\subsection{2. Methods}\label{methods}}

\hypertarget{network-analysis}{%
\subsubsection{2.1 Network analysis}\label{network-analysis}}

\emph{Psychological networks} have their own characteristics which
differentiate them from other networks of the graph theory (Epskamp et
al., 2017; Newman, 2010). Nodes are normally some measures or items of a
psychological assessment tool, the edges of the network reflect their
estimates of association derived from the data.

\hypertarget{removing-redundant-items}{%
\paragraph{2.1.1 Removing redundant
items}\label{removing-redundant-items}}

As a preparation step, we applied the \emph{goldbricker} function from
the R package \emph{networktools} to search for pairs of items in the
RRS, SCRS and MHS items which were highly correlated (r \textgreater{}
0.50) and showed highly the same pattern with other items. When such
``bad pairs'' were identified, reduction techniques were put into use
(with the \emph{net\_reduce} function in R): instead of the pair, the
first principal component of the two variables were kept in the data as
a new, merged variable.

The RRS items required no reduction steps as the highly correlating
items did not extend their pattern to more nodes (their correlations to
other nodes did not violate the criteria that \textgreater{} 75\% of
these correlations did not differ significantly).

\hypertarget{network-analysis-with-lasso}{%
\paragraph{2.1.2 Network Analysis with
LASSO}\label{network-analysis-with-lasso}}

For the network analysis, the tuning parameter of \emph{gamma} in the
grahical LASSO alogrithm has been set to 0.5 value in accordance with
(Bernstein et al., 2019) in order to achieve sparse graphs. The final
networks were chosen by the lowest values of the \emph{extended Bayesian
Information Criterion} (EBIC).

\hypertarget{community-detection}{%
\paragraph{2.1.3. Community detection}\label{community-detection}}

In order to get an understanding of community of nodes and to see if
nodes can be grouped into neighbourhoods, we used the \emph{igraph} R
package. When detecting possible communities of nodes, the detection
happened with the following parameters in the \emph{spinglass
algorithm}: TODO: try the walktrap algorithm from igraph!

\hypertarget{expected-influence-of-nodes}{%
\paragraph{2.1.4. Expected Influence of
nodes}\label{expected-influence-of-nodes}}

To understand certain items' centrality and importance, besides
regularly used centrality measures (such as strength, closeness and
betweenness of nodes in a network), one can also use the One- and
Two-Step Expected Influence (EI) measures (Robinaugh et al, 2016).
Traditionally in unweighted networks (where network edges are encoded to
0 or 1) a node's centrality is the sum of edges, and in weighted nodes
it is the sum of the edge weights. In this approach negative edges are
treated the same way as positives, as the absolute values are used for
centrality estimation. On the other hand, the Expected Influence metrics
take into account if a node has positive and negative edges, too.

\hypertarget{identifying-important-nodes}{%
\paragraph{2.1.5. Identifying important
nodes}\label{identifying-important-nodes}}

In order to further elaborate the question of importance within the
network, we identified \emph{bridge nodes} using the functions provided
by the \emph{networktools} R package (Jones, 2017). We estimated both
one-step and two-step bridge expected influence measures which reflect
the sum of edge weights connecting one node to other groups on nodes (or
communities of nodes); and also express a node's secondary influence on
the rest of the network, respectively. Higher influence scores indicate
that a certain node is more likely to active neighbouring nodes,
subgroups of nodes or communities.

\hypertarget{results}{%
\subsection{3. Results}\label{results}}

\hypertarget{network-plots}{%
\subsubsection{3.1 Network plots}\label{network-plots}}

RSS network

\includegraphics{main_files/figure-latex/unnamed-chunk-4-1.pdf}

SCRS network
\includegraphics{main_files/figure-latex/unnamed-chunk-5-1.pdf}

\begin{verbatim}
##   |                                                                              |                                                                      |   0%  |                                                                              |====                                                                  |   5%  |                                                                              |=======                                                               |  10%  |                                                                              |==========                                                            |  15%  |                                                                              |==============                                                        |  20%  |                                                                              |==================                                                    |  25%  |                                                                              |=====================                                                 |  30%  |                                                                              |========================                                              |  35%  |                                                                              |============================                                          |  40%  |                                                                              |================================                                      |  45%  |                                                                              |===================================                                   |  50%  |                                                                              |======================================                                |  55%  |                                                                              |==========================================                            |  60%  |                                                                              |==============================================                        |  65%  |                                                                              |=================================================                     |  70%  |                                                                              |====================================================                  |  75%  |                                                                              |========================================================              |  80%  |                                                                              |============================================================          |  85%  |                                                                              |===============================================================       |  90%  |                                                                              |==================================================================    |  95%  |                                                                              |======================================================================| 100%
\end{verbatim}

\begin{verbatim}
## SCRS_4 & SCRS_2 SCRS_8 & SCRS_1 
##       0.1666667       0.2222222
\end{verbatim}

\includegraphics{main_files/figure-latex/unnamed-chunk-6-1.pdf}

\hypertarget{community-detection-1}{%
\subsubsection{3.2 Community detection}\label{community-detection-1}}

Community mapping of RSS, MHC and SCRS with the \emph{spinglass
algorithm}:

RSS communities identified with parameters \emph{gamma = 0.5, spins =
25, starting temperature = 1, ending temperature = 0.01, cooling factor
= 0.99}.

\includegraphics{main_files/figure-latex/unnamed-chunk-7-1.pdf}

SCRS communities identified with parameters \emph{gamma = 0.5, spins =
25, starting temperature = 1, ending temperature = 0.01, cooling factor
= 0.99}.

\includegraphics{main_files/figure-latex/unnamed-chunk-8-1.pdf}

RRS and MHC combined communities identified with parameters \emph{gamma
= 0.5, spins = 25, starting temperature = 1, ending temperature = 0.01,
cooling factor = 0.99}.

\includegraphics{main_files/figure-latex/unnamed-chunk-10-1.pdf}

\hypertarget{expected-influence}{%
\subsubsection{3.3 Expected Influence}\label{expected-influence}}

In order to estimate the Expected Influence metrics, we used the
\emph{expectedInf} function of the \emph{networktools} R package. The
influence score are normalized.

RRS Expected Influence Measures

\begin{verbatim}
## [1] "RRS"
## [1] "RRS"
## [1] "RRS"
## [1] "RRS"
## [1] "RRS"
## [1] "RRS"
## [1] "RRS"
## [1] "RRS"
## [1] "RRS"
## [1] "RRS"
\end{verbatim}

\includegraphics{main_files/figure-latex/unnamed-chunk-11-1.pdf}

SCRS Expected Influence Measures

\begin{verbatim}
## [1] "SCRS"
## [1] "SCRS"
## [1] "SCRS"
## [1] "SCRS"
## [1] "SCRS"
## [1] "SCRS"
## [1] "SCRS"
## [1] "SCRS"
## [1] "SCRS"
## [1] "SCRS"
\end{verbatim}

\includegraphics{main_files/figure-latex/unnamed-chunk-13-1.pdf}

Combined RRS and SCRS Expected Influence Measures

\begin{verbatim}
## [1] "RRS"
## [1] "RRS"
## [1] "RRS"
## [1] "RRS"
## [1] "RRS"
## [1] "RRS"
## [1] "RRS"
## [1] "RRS"
## [1] "RRS"
## [1] "RRS"
## [1] "SCRS"
## [1] "SCRS"
## [1] "SCRS"
## [1] "SCRS"
## [1] "SCRS"
## [1] "SCRS"
## [1] "SCRS"
## [1] "SCRS"
\end{verbatim}

\includegraphics{main_files/figure-latex/unnamed-chunk-14-1.pdf}

\hypertarget{bridge-nodes}{%
\subsubsection{3.4 Bridge nodes}\label{bridge-nodes}}

Using the \emph{networktools} R library bridge expected influence
measure were estimated for the three tools we used in this study. We
also added aesthetics to reflect which items belonged to the same
community according to our previous clustering steps.

RRS bridge estimates

\includegraphics{main_files/figure-latex/unnamed-chunk-15-1.pdf}

SCRS bridge estimates

\includegraphics{main_files/figure-latex/unnamed-chunk-16-1.pdf}

Combined RRS and SCRS bridge estimates

\includegraphics{main_files/figure-latex/unnamed-chunk-18-1.pdf}

\hypertarget{references}{%
\subsection*{References}\label{references}}
\addcontentsline{toc}{subsection}{References}

\hypertarget{refs}{}
\leavevmode\hypertarget{ref-Bernstein2019}{}%
Bernstein, E. E., Heeren, A., \& McNally, R. J. (2019). Reexamining
trait rumination as a system of repetitive negative thoughts: A network
analysis. \emph{Journal of Behavior Therapy and Experimental
Psychiatry}, \emph{63}, 21--27.
\url{https://doi.org/10.1016/j.jbtep.2018.12.005}

\leavevmode\hypertarget{ref-Epskamp2017}{}%
Epskamp, S., Borsboom, D., \& Fried, E. I. (2017). Estimating
psychological networks and their accuracy: A tutorial paper.
\emph{Behavior Research Methods}, \emph{50}(1), 195--212.
\url{https://doi.org/10.3758/s13428-017-0862-1}

\leavevmode\hypertarget{ref-Jones2017}{}%
Jones, P. (2017). \emph{Networktools: Assorted and tools for identifying
and important nodes and in networks}.

\leavevmode\hypertarget{ref-Newman2010}{}%
Newman, M. E. J. (2010). \emph{Networks - an introduction}.

\end{document}
